\chapter{Electrochemical Systems}\label{eq:electrochemistry}
Electrochemists study chemical reactions that involve electron
transfer for many reasons such as biochemical sensing,
electroplating, energy storage, imaging etc. A popular technique to
understand the diffusive or absortive nature, thermodynamic
parameters, as well as the existence of coupled homogenous chemical
reactions is cyclic voltammetry.

\section{Mathematical model of an electrochemical system}

We can represent a general reduction-oxidation (redox) process at the
electrode as:
\begin{equation}\label{eq:redox}
  \mathrm{A} + \mathrm{e}^{-} \rightleftharpoons \mathrm{B}
\end{equation}
In a typical experiment, the potential is swept from some starting
potential, $\mathrm{E}_i$ where the species A is not electroreduced,
to some -- more negative -- vertex potential, $\mathrm{E}_v$, at
which the electron transfer between the species A and the electrode
is rapid and the species B is formed. The potential is then swept
back to $E_i$ causing electron transfer in the opposite direction and
thus the reformation of A. Throughout this process the current, $I$,
is recorded. Plotting the current against the potential gives the
characteristic cyclic voltammagram in

\begin{equation}\label{eq:bv_equation}
  \frac{\partial \mathrm{C}(X, T)}{\partial T} = \frac{\partial^2
  C(X, T)}{\partial X^2}
\end{equation}
Subject to:
\begin{equation}
  \begin{split}
    C(X,0) &= 1 \\
    \left.\frac{\partial C(X,T)}{\partial X}\right|_{X=0}
    &= K_0\!\left( C(0,T)e^{-\alpha\theta} -
    \big(1-C(0,T)\big)e^{(1-\alpha)\theta}\right)
  \end{split}
\end{equation}
